\section{Comparisons, variations and supporting calculations}
\label{app:coupling}
The calculations here support the field comparison, phase connection and coupled response in Sections~\ref{sec:field}--\ref{sec:color}. First, explicit matrix examples verify the comparisons used in the construction. The later derivatives follow a source variation through its responses and compatibility conditions. A curve parameter $s$ labels these variations and is independent of physical time.

\subsection{Matrix comparisons}
For the two matrices used to explain the field pairing, transposition exchanges rows and columns,
\FieldExampleTransposeDisplay
Multiplication then collects matching column products on the diagonal,
\FieldExampleDiagonalDisplay
and their sum gives
\FieldExampleTraceDisplay
This verifies the trace formula used in Section~\ref{sec:field}.

A self-adjoint map transfers its action between the entries of the inner product. For example,
\FieldSelfAdjointExampleDisplay
gives
\FieldSelfAdjointActionsDisplay
and hence
\FieldSelfAdjointComparisonsDisplay
The equality holds for every vector pair because the matrix is symmetric,
\FieldSelfAdjointIdentityDisplay
This identifies symmetric matrices as the self-adjoint component values used by the field.

The product metric applies the same comparison at every component label. Consider two complete field changes with components
\FieldProductFirstExampleDisplay
and
\FieldProductSecondExampleDisplay
Their component pairings are
\FieldProductComparisonsDisplay
so the complete comparison is
\FieldProductTotalDisplay
Self-pairing similarly gives
\FieldProductLengthDisplay
Each labelled change contributes to the length of the complete variation.

\subsection{Phase rotation and moving frames}
In the oriented plane of Section~\ref{sec:phase}, rotation through $\FieldPlaneAngle$ multiplies its complex coordinate by
\FieldPlanePhaseFactorDisplay
The real coordinates consequently transform as
\FieldPlaneRotationCoordinatesDisplay
Since
\FieldPlanePhaseLengthDisplay
the rotation preserves length. These factors constitute
\FieldPlanePhaseGroupDisplay
and compose by addition of angles,
\FieldPlanePhaseCompositionDisplay
This is the phase action carried by the radial and transverse plane.

To verify the connection's role, keep a plane fixed and rotate its frame,
\FieldConnectionRotatingFrameDisplay
Its angular rate is
\FieldConnectionRotatingRateDisplay
Now take the fixed vector
\FieldConnectionFixedVectorDisplay
Its coordinates in that rotating frame are
\FieldConnectionFixedCoordinatesDisplay
or, in complex form,
\FieldConnectionFixedComplexDisplay
Their ordinary derivative is nonzero, but the frame term cancels it,
\FieldConnectionCancellationDisplay
The covariant derivative therefore records the fixed geometric vector correctly.

An additional frame angle $\FieldGaugeAngle$ changes the basis by
\FieldConnectionFrameChangeDisplay
The coordinate of the same vector changes oppositely,
\FieldConnectionCoordinateChangeDisplay
while the connection becomes
\FieldConnectionPotentialChangeDisplay
Together these transformations give
\FieldConnectionCovarianceDisplay
so either frame describes the same derivative. Choosing $\FieldGaugeUnrotation$ removes the earlier frame rotation and gives constant coordinates for the fixed vector.

\subsection{Weak waves and collective mass}
For the conventional gluon connection, linearization about the trivial connection gives the transverse wave equation
\FieldYMWaveDisplay
Here $z$ denotes a spatial coordinate. With angular frequency $\FieldYMFrequency$ and wave vector $\FieldYMWaveVector$, it implies
\FieldYMDispersionDisplay
Using $\FieldYMQuantumEnergy$ and $\FieldYMQuantumMomentum$, where $\hbar$ is the reduced Planck constant, gives
\FieldYMMasslessEnergyDisplay
Comparison with
\FieldYMRelativisticEnergyDisplay
identifies $m=0$ for these weak-field excitations~\cite{jaffeWitten}.

Their zero individual masses do not require zero mass for a collection. If two excitations have energy $\FieldYMExampleEnergy$ and opposite momenta,
\FieldYMOppositeMomentaDisplay
the totals are
\FieldYMCombinedEnergyDisplay
and their invariant mass is
\FieldYMPairMassDisplay
This illustrates the distinction between constituent mass and collective mass; binding remains a separate dynamical property.

In pure gauge theory, a glueball is a bound excitation of the interacting gluon field. Lattice calculations estimate such collective energies. Morningstar and Peardon obtained approximately
\FieldYMGlueballMassDisplay
for the lightest scalar state, using their stated conversion to physical units~\cite{morningstarPeardon1999}. The label $\FieldYMScalarLabel$ denotes zero spin and positive parity and charge-conjugation parity. This result is a comparison, not an input to the constructed source's mass.

A common mass gap would require a positive threshold for all physical excitations. The opposite-momentum example alone supplies no such threshold because its input energy can be arbitrarily small. The main construction instead follows the specific closed geometry to its complete energy--momentum.

\subsection{Derivatives of the coupling}
The response used in Section~\ref{sec:coupling} is
\FieldCouplingResponseDisplay
with component rule
\FieldCouplingWeightsDisplay
Both factors depend on the source. Writing
\FieldCouplingDerivativeDisplay
gives
\FieldCouplingVariationDisplay
The two terms vary the component and its pairings with the other components.

Writing the same response directly as a polynomial in its input gives
\CouplingMapPolynomialCouplingDisplay
Its differential is the preceding variation with $L_X=D\Lambda(X)$ written out,
\CouplingMapCouplingDifferentialDisplay
The derivative of its scalar coefficient is
\DifferentiationCoefficientDisplay
Each pairing varies through both entries, so an unchanged input component can still acquire a changed response.

Differentiating once more in independent directions $\delta_1X,\delta_2X$ gives
\CouplingMapCouplingSecondDifferentialDisplay
The first two terms differentiate the two factors of the first variation; the last differentiates the pairings again. Exchange of the variation directions leaves the result unchanged.

For the matrix example in Section~\ref{sec:coupling}, let the first component vary in its second diagonal entry,
\DifferentiationMatrixCurveDisplay
The coefficient $w_1=w_2=2+s$ varies with that entry, giving at $s=0$
\DifferentiationMatrixDerivativesDisplay
The quadratic term comes from their product. This verifies why coefficient variation must remain in the response derivative.

\subsection{Successive responses}
To carry the variation into the next level, write
\DifferentiationChainDefinitionDisplay
The derivatives $L_X$ and $L_Y$ are evaluated at different inputs. Since the second input is itself a response, its change must also be differentiated.

For a configuration direction $\DifferentiationBaseDirection$, put $\delta X=DX[\eta]$. The chain rule gives
\CouplingMapFullChainRuleDisplay
The second response receives the composed change $L_YL_X\delta X$, while the last expression differentiates the evaluation point of $L_Y$.

If the varied vector also depends on the configuration, the full derivative is
\CouplingMapBothDifferentialVariationsDisplay
All three contributions enter the changing frame, along with derivatives of its projections and normalization. Treating either response derivative as fixed would omit part of that change.

\subsection{Normalization and compatibility}
Normalization itself contributes nonlinear derivatives. For example,
\DifferentiationNormalizationDisplay
The unnormalized vector has zero second derivative, but its unit direction does not. Differentiating unit length in general gives
\DifferentiationUnitConstraintDisplay
The first identity makes velocity tangent to the unit sphere; the second fixes the acceleration's normal component.

Compatibility must be differentiated along the same source curve. Let $\CombinedFieldMap$ collect the responses, metric, connection and derivatives in a smooth condition $\CompatibilityEquality$. Take a twice differentiable curve $\CompatibleFieldCurve$ satisfying it, and set $\FieldVelocityAcceleration$. Dots refer to $s$, with derivatives of $\Phi$ evaluated at $X(0)$ and those of $f$ at $\Phi(X(0))$.

One and two differentiations give
\CouplingMapPropagatedCompatibilityDisplay
The second-order condition contains source acceleration, nonlinear response and curvature of the constraint. Retaining all three terms ensures that the varied induced quantities follow the source compatibility.

For a smooth inequality $\CompatibilityInequality$ that is an equality initially, a compatible curve with $s\ge0$ obeys
\CouplingMapContactDerivativeDisplay
Writing $\DifferentiationContactFunction$ gives the Taylor expansion
\DifferentiationContactTaylorDisplay
A negative first derivative violates the inequality for sufficiently small positive $s$. If the first derivative vanishes, the second must be nonnegative. A curve allowed for both signs requires zero first derivative.

The distinction is illustrated by
\DifferentiationContactExamplesDisplay
The linear example is admissible on one side and the quadratic example on both. These differential restrictions are necessary; the curve must still satisfy the original compatibility relations.

\subsection{Graph and normal planes}
The first response places a variation in the graph
\FieldCouplingGraphDisplay
The adjoint relation
\FieldCouplingAdjointDisplay
then gives a normal lift
\FieldCouplingNormalDisplay
whose complete pairing with the graph is
\FieldCouplingOrthogonalityDisplay
The cancellation establishes the perpendicular directions used in Section~\ref{sec:coupling}.

Inserting the radial and transverse frame in the selected component gives
\FieldCouplingGraphPlaneDisplay
and
\FieldCouplingNormalPlaneDisplay
The graph retains the two independent directions in its first position; the normal lift retains them in its second. Their ranks therefore remain two.

For an explicit check, take
\FieldCouplingExampleDerivativeDisplay
The paired directions are
\FieldCouplingExampleLiftsDisplay
and, for example,
\FieldCouplingExampleCancellationDisplay
All cross-plane pairings vanish similarly. Each vector has squared length five, so division by $\FieldCouplingExampleNormalization$ gives an orthonormal frame within either plane.

For the three-level joint variation of Section~\ref{sec:color}, the lifts become
\AlgebraTripleLiftDisplay
The continuing matrix example gives
\AlgebraResponseExampleDisplay
for the final graph component. Adjacent adjoint cancellations show
\AlgebraNormalCancellationDisplay
Both normal planes are thus perpendicular to the graph, though they may overlap each other through the middle position.

Subtracting that overlap gives
\AlgebraOrthogonalColumnsDisplay
Here $\mathsf W_2$ is the second normal map before projection. For a plane-coordinate column $\xi\in\mathbb R^2$, the coefficients $\alpha\in\mathbb R^2$ solve
\AlgebraProjectionDisplay
The inverse exists because $\mathsf V_1$ has independent columns. Subtraction leaves the third block of $\mathsf V_2$ equal to $\mathsf V$, since $\mathsf V_1$ has zero third block. All three maps therefore retain two independent columns and the mutual orthogonality used in the main text.

\subsection{Pair-generator calculations}
The matrix units act by
\AlgebraMatrixUnitExampleDisplay
so their product connects the first and third coordinates. For the real pair generator,
\AlgebraGeneratorActionDisplay
multiplication in both orders gives
\AlgebraBracketProductsDisplay
This verifies the first--third generator in Section~\ref{sec:color}. The imaginary generator follows from the same matrix-unit product.

For the diagonal components, the trace pairing gives
\AlgebraDiagonalNormExampleDisplay
and the first--third difference is
\AlgebraDiagonalExpansionDisplay
These relations fix their normalization and dependence.

Finally, connection curvature retains the order of its matrix-valued factors. For local tangent directions $\eta,\zeta$,
\AlgebraWedgeExampleDisplay
The nonzero example illustrates the adjacent interaction in the first--third curvature entry.
